\documentclass[10pt,a4paper]{amsart}
\usepackage{amsmath,amssymb,amsthm,graphicx}
% Retained preamble for the Stacks Project chapter snapshot in this directory.
% The chapter's own preamble.tex is a copy of the shared Stacks Project preamble; the block
% below is sliced MECHANICALLY, line for line, from it. Nothing was added, removed or
% redefined; the wrapper class, the xr-hyper/hyperref wiring and the externaldocument list
% are omitted because the assembler supplies the document class.
%
% lines taken from preamble.tex: 3, 5, 9, 11, 12, 156, 157, 158, ... (51 lines total)
%
% package coverage restored from the chapter's own preamble (2026-10-13): xy, xy-2cell

\usepackage[all]{xy}
\xyoption{2cell}
\UseAllTwocells

\newenvironment{reference}{\comment}{\endcomment}
\newenvironment{slogan}{\comment}{\endcomment}
\newenvironment{history}{\comment}{\endcomment}
\theoremstyle{plain}
\newtheorem{theorem}[subsection]{Theorem}
\newtheorem{proposition}[subsection]{Proposition}
\newtheorem{lemma}[subsection]{Lemma}
\theoremstyle{definition}
\newtheorem{definition}[subsection]{Definition}
\newtheorem{example}[subsection]{Example}
\newtheorem{exercise}[subsection]{Exercise}
\newtheorem{situation}[subsection]{Situation}
\theoremstyle{remark}
\newtheorem{remark}[subsection]{Remark}
\newtheorem{remarks}[subsection]{Remarks}
\numberwithin{equation}{subsection}
\def\lim{\mathop{\mathrm{lim}}\nolimits}
\def\colim{\mathop{\mathrm{colim}}\nolimits}
\def\Spec{\mathop{\mathrm{Spec}}}
\def\Hom{\mathop{\mathrm{Hom}}\nolimits}
\def\Ext{\mathop{\mathrm{Ext}}\nolimits}
\def\SheafHom{\mathop{\mathcal{H}\!\mathit{om}}\nolimits}
\def\SheafExt{\mathop{\mathcal{E}\!\mathit{xt}}\nolimits}
\def\Sch{\mathit{Sch}}
\def\Mor{\mathop{\mathrm{Mor}}\nolimits}
\def\Ob{\mathop{\mathrm{Ob}}\nolimits}
\def\Sh{\mathop{\mathit{Sh}}\nolimits}
\def\NL{\mathop{N\!L}\nolimits}
\def\CH{\mathop{\mathrm{CH}}\nolimits}
\def\proetale{{pro\text{-}\acute{e}tale}}
\def\etale{{\acute{e}tale}}
\def\QCoh{\mathit{QCoh}}
\def\Ker{\mathop{\mathrm{Ker}}}
\def\Im{\mathop{\mathrm{Im}}}
\def\Coker{\mathop{\mathrm{Coker}}}
\def\Coim{\mathop{\mathrm{Coim}}}
\DeclareMathSymbol{\boxtimes}{\mathbin}{AMSa}{"02}
\def\QCohstack{\mathcal{QC}\!\mathit{oh}}
\def\Cohstack{\mathcal{C}\!\mathit{oh}}
\def\Spacesstack{\mathcal{S}\!\mathit{paces}}
\def\Quotfunctor{\mathrm{Quot}}
\def\Hilbfunctor{\mathrm{Hilb}}
\def\Curvesstack{\mathcal{C}\!\mathit{urves}}
\def\Polarizedstack{\mathcal{P}\!\mathit{olarized}}
\def\Complexesstack{\mathcal{C}\!\mathit{omplexes}}
\def\Pic{\mathop{\mathrm{Pic}}\nolimits}
\def\Picardstack{\mathcal{P}\!\mathit{ic}}
\def\Picardfunctor{\mathrm{Pic}}
\def\Deformationcategory{\mathcal{D}\!\mathit{ef}}
\usepackage{cleveref}
\crefname{originallemma}{original lemma}{original lemmas}
\Crefname{originallemma}{Original lemma}{Original lemmas}
\makeatletter
\newlabel{algebra-lemma-NL-homotopy}{{00S1}{source}{}{sourcepointer-1}{}}
\newlabel{algebra-lemma-NL-homotopy@cref}{{[originallemma][0][]00S1}{[1][0][]source}{}{}{}}
\makeatother
\crefname{originallemma}{original lemma}{original lemmas}
\Crefname{originallemma}{Original lemma}{Original lemmas}
\makeatletter
\newlabel{algebra-lemma-complete-modulo-nilpotent}{{0DYC}{source}{}{sourcepointer-2}{}}
\newlabel{algebra-lemma-complete-modulo-nilpotent@cref}{{[originallemma][0][]0DYC}{[1][0][]source}{}{}{}}
\makeatother
\crefname{originallemma}{original lemma}{original lemmas}
\Crefname{originallemma}{Original lemma}{Original lemmas}
\makeatletter
\newlabel{algebra-lemma-completion-tensor}{{00MA}{source}{}{sourcepointer-3}{}}
\newlabel{algebra-lemma-completion-tensor@cref}{{[originallemma][0][]00MA}{[1][0][]source}{}{}{}}
\makeatother
\crefname{originallemma}{original lemma}{original lemmas}
\Crefname{originallemma}{Original lemma}{Original lemmas}
\makeatletter
\newlabel{algebra-section-netherlander}{{00S0}{source}{}{sourcepointer-4}{}}
\newlabel{algebra-section-netherlander@cref}{{[originallemma][0][]00S0}{[1][0][]source}{}{}{}}
\makeatother
\crefname{originallemma}{original lemma}{original lemmas}
\Crefname{originallemma}{Original lemma}{Original lemmas}
\makeatletter
\newlabel{lemma-completion}{{0DYV}{source}{}{sourcepointer-5}{}}
\newlabel{lemma-completion@cref}{{[originallemma][0][]0DYV}{[1][0][]source}{}{}{}}
\makeatother
\crefname{originallemma}{original lemma}{original lemmas}
\Crefname{originallemma}{Original lemma}{Original lemmas}
\makeatletter
\newlabel{more-algebra-lemma-ext-annihilated-into}{{0DYJ}{source}{}{sourcepointer-6}{}}
\newlabel{more-algebra-lemma-ext-annihilated-into@cref}{{[originallemma][0][]0DYJ}{[1][0][]source}{}{}{}}
\makeatother
\crefname{originallemma}{original lemma}{original lemmas}
\Crefname{originallemma}{Original lemma}{Original lemmas}
\makeatletter
\newlabel{more-algebra-lemma-neighbourhood-equivalence}{{05EC}{source}{}{sourcepointer-7}{}}
\newlabel{more-algebra-lemma-neighbourhood-equivalence@cref}{{[originallemma][0][]05EC}{[1][0][]source}{}{}{}}
\makeatother
\crefname{originallemma}{original lemma}{original lemmas}
\Crefname{originallemma}{Original lemma}{Original lemmas}
\makeatletter
\newlabel{more-algebra-lemma-pseudo-coherence-and-ext}{{087Q}{source}{}{sourcepointer-8}{}}
\newlabel{more-algebra-lemma-pseudo-coherence-and-ext@cref}{{[originallemma][0][]087Q}{[1][0][]source}{}{}{}}
\makeatother
\title{Stacks, Examples of Deformations: lemma-first-order-completion}
\author{The Stacks Project Authors}
\date{stacks-project:examples-defos.tex:local-snapshot-20261012; GNU FDL 1.2 or later}
\begin{document}
\maketitle
\noindent\textit{Editorial note.} Statement span is lines 2382-2393 of sources/stacks-examples-defos/pinned-original-source.tex; proof span is lines 2395-2489, the complete author proof (93 body lines).\par\medskip
\noindent\textit{Editorial note.} Navigation stubs: 8 cited-but-not-included labels are supplied as non-load-bearing pointers with their pinned tags where the snapshot tag index has them.\par\medskip
\noindent\textit{Editorial note.} extra\_packages is empty on purpose: the retained preamble already loads xy with options and a second declaration clashes in the assembler wrapper.\par\medskip
\noindent\textit{Editorial note.} Text is retained verbatim, including the author's own typos and placeholders.\par\medskip
\noindent\textit{Editorial note.} Queue index 0 of queue\_examples-defos.json (pb=3523, stmt 2382-2393, proof 2395-2489).\par\medskip
\noindent\textit{Allowed standard prerequisites (editorial).}\begin{itemize}
\item Explicit prerequisite boundary. The cited inputs are prerequisites; the argument in the retained span is not replaced by any cited result.
\end{itemize}\medskip
\ifdefined\hypertarget\hypertarget{sourcepointer-1}{}\fi
\noindent\textit{Source reference.} Original lemma 00S1: sources/stacks-examples-defos/pinned-original-source.tex, source lines 2489--2489. Cited by the retained statement or proof but not part of the collected spans; supplied only as a navigation pointer. The argument does not depend on its content.\par\medskip
\ifdefined\hypertarget\hypertarget{sourcepointer-2}{}\fi
\noindent\textit{Source reference.} Original lemma 0DYC: sources/stacks-examples-defos/pinned-original-source.tex, source lines 2489--2489. Cited by the retained statement or proof but not part of the collected spans; supplied only as a navigation pointer. The argument does not depend on its content.\par\medskip
\ifdefined\hypertarget\hypertarget{sourcepointer-3}{}\fi
\noindent\textit{Source reference.} Original lemma 00MA: sources/stacks-examples-defos/pinned-original-source.tex, source lines 2489--2489. Cited by the retained statement or proof but not part of the collected spans; supplied only as a navigation pointer. The argument does not depend on its content.\par\medskip
\ifdefined\hypertarget\hypertarget{sourcepointer-4}{}\fi
\noindent\textit{Source reference.} Original lemma 00S0: sources/stacks-examples-defos/pinned-original-source.tex, source lines 2489--2489. Cited by the retained statement or proof but not part of the collected spans; supplied only as a navigation pointer. The argument does not depend on its content.\par\medskip
\ifdefined\hypertarget\hypertarget{sourcepointer-5}{}\fi
\noindent\textit{Source reference.} Original lemma 0DYV: sources/stacks-examples-defos/pinned-original-source.tex, source lines 1309--1318. Cited by the retained statement or proof but not part of the collected spans; supplied only as a navigation pointer. The argument does not depend on its content.\par\medskip
\ifdefined\hypertarget\hypertarget{sourcepointer-6}{}\fi
\noindent\textit{Source reference.} Original lemma 0DYJ: sources/stacks-examples-defos/pinned-original-source.tex, source lines 2489--2489. Cited by the retained statement or proof but not part of the collected spans; supplied only as a navigation pointer. The argument does not depend on its content.\par\medskip
\ifdefined\hypertarget\hypertarget{sourcepointer-7}{}\fi
\noindent\textit{Source reference.} Original lemma 05EC: sources/stacks-examples-defos/pinned-original-source.tex, source lines 2489--2489. Cited by the retained statement or proof but not part of the collected spans; supplied only as a navigation pointer. The argument does not depend on its content.\par\medskip
\ifdefined\hypertarget\hypertarget{sourcepointer-8}{}\fi
\noindent\textit{Source reference.} Original lemma 087Q: sources/stacks-examples-defos/pinned-original-source.tex, source lines 2489--2489. Cited by the retained statement or proof but not part of the collected spans; supplied only as a navigation pointer. The argument does not depend on its content.\par\medskip
\begin{lemma}
\label{lemma-first-order-completion}
Let $k$ be a field. Let $B$ be a finite type $k$-algebra.
Let $J \subset B$ be an ideal such that
$\Spec(B) \to \Spec(k)$ is smooth on the complement of $V(J)$.
Let $N$ be a finite $B$-module.
Then there is a canonical bijection
$$
\text{Exal}_k(B, N) \to \text{Exal}_k(B^\wedge, N^\wedge)
$$
Here $B^\wedge$ and $N^\wedge$ are the $J$-adic completions.
\end{lemma}

\begin{proof}
The map is given by completion: given $0 \to N \to C \to B \to 0$
in $\text{Exal}_k(B, N)$ we send it to the completion $C^\wedge$
of $C$ with respect to the inverse image of $J$. Compare with
the proof of Lemma \ref{lemma-completion}.

\medskip\noindent
Since $k \to B$ is of finite presentation the complex
$\NL_{B/k}$ can be represented by a complex
$N^{-1} \to N^0$ where $N^i$ is a finite $B$-module, see
Algebra, Section \ref{algebra-section-netherlander} and
in particular
Algebra, Lemma \ref{algebra-lemma-NL-homotopy}.
As $B$ is Noetherian, this means that $\NL_{B/k}$
is pseudo-coherent. For $g \in J$ the $k$-algebra $B_g$
is smooth and hence $(\NL_{B/k})_g = \NL_{B_g/k}$
is quasi-isomorphic to a finite projective $B$-module sitting in degree $0$.
Thus $\text{Ext}^i_B(\NL_{B/k}, N)_g = 0$ for $i \geq 1$
and any $B$-module $N$. By
More on Algebra, Lemma \ref{more-algebra-lemma-ext-annihilated-into}
we conclude that
$$
\text{Ext}^1_B(\NL_{B/k}, N) \longrightarrow
\lim_n \text{Ext}^1_B(\NL_{B/k}, N/J^n N)
$$
is an isomorphism for any finite $B$-module $N$.

\medskip\noindent
Injectivity of the map.
Suppose that $0 \to N \to C \to B \to 0$ is in $\text{Exal}_k(B, N)$
and maps to zero in $\text{Exal}_k(B^\wedge, N^\wedge)$.
Choose a splitting $C^\wedge = B^\wedge \oplus N^\wedge$.
Then the induced map $C \to C^\wedge \to N^\wedge$
gives maps $C \to N/J^nN$ for all $n$.
Hence we see that our element is in the kernel of the maps
$$
\text{Ext}^1_B(\NL_{B/k}, N) \to
\text{Ext}^1_B(\NL_{B/k}, N/J^n N)
$$
for all $n$. By the previous paragraph we conclude that
our element is zero.

\medskip\noindent
Surjectivity of the map. Let $0 \to N^\wedge \to C' \to B^\wedge \to 0$
be an element of $\text{Exal}_k(B^\wedge, N^\wedge)$.
Pulling back by $B \to B^\wedge$ we get an element
$0 \to N^\wedge \to C'' \to B \to 0$ in
$\text{Exal}_k(B, N^\wedge)$.
we have
$$
\text{Ext}^1_B(\NL_{B/k}, N^\wedge) =
\text{Ext}^1_B(\NL_{B/k}, N) \otimes_B B^\wedge =
\text{Ext}^1_B(\NL_{B/k}, N)
$$
The first equality as $N^\wedge = N \otimes_B B^\wedge$
(Algebra, Lemma \ref{algebra-lemma-completion-tensor})
and
More on Algebra, Lemma \ref{more-algebra-lemma-pseudo-coherence-and-ext}.
The second equality because $\text{Ext}^1_B(\NL_{B/k}, N)$
is $J$-power torsion (see above), $B \to B^\wedge$ is flat and induces
an isomorphism $B/J \to B^\wedge/JB^\wedge$, and
More on Algebra, Lemma \ref{more-algebra-lemma-neighbourhood-equivalence}.
Thus we can find a $C \in \text{Exal}_k(B, N)$ mapping to $C''$ in
$\text{Exal}_k(B, N^\wedge)$.
Thus
$$
0 \to N^\wedge \to C' \to B^\wedge \to 0
\quad\text{and}\quad
0 \to N^\wedge \to C^\wedge \to B^\wedge \to 0
$$
are two elements of $\text{Exal}_k(B^\wedge, N^\wedge)$
mapping to the same element of $\text{Exal}_k(B, N^\wedge)$.
Taking the difference we get an element
$0 \to N^\wedge \to C' \to B^\wedge \to 0$ of
$\text{Exal}_k(B^\wedge, N^\wedge)$
whose image in $\text{Exal}_k(B, N^\wedge)$ is zero.
This means there exists
$$
\xymatrix{
0 \ar[r] &
N^\wedge \ar[r] &
C' \ar[r] &
B^\wedge \ar[r] & 0 \\
& & B \ar[u]^\sigma \ar[ru]
}
$$
Let $J' \subset C'$ be the inverse image of $JB^\wedge \subset B^\wedge$.
To finish the proof it suffices to note that
$\sigma$ is continuous for the $J$-adic topology on $B$
and the $J'$-adic topology on $C'$ and that $C'$ is $J'$-adically complete by
Algebra, Lemma \ref{algebra-lemma-complete-modulo-nilpotent}
(here we also use that $C'$ is Noetherian; small detail omitted).
Namely, this means that $\sigma$ factors through the
completion $B^\wedge$ and $C' = 0$ in $\text{Exal}_k(B^\wedge, N^\wedge)$.
\end{proof}

\end{document}
